% exp-fips203-mlkem-pke-alg13-16171820-2s1e-k4 — Alg.13 行 16–20（2s1e UB 融合，FIPS 合规采样）
% 编译：bash ../../../../scripts/xelatex-clean.sh exp-fips203-mlkem-pke-alg13-16171820-2s1e-k4-实现方案-customspec.tex
\documentclass[UTF8,a4paper,11pt]{ctexart}
\usepackage{amsmath,amssymb}
\usepackage{booktabs}
\usepackage{geometry}
\usepackage{hyperref}
\usepackage{longtable}
\usepackage{array}
\usepackage{seqsplit}
\geometry{margin=2.2cm}
\newcolumntype{L}[1]{>{\raggedright\arraybackslash}p{#1}}
\newcommand{\apiname}[1]{\texttt{\seqsplit{#1}}}
\newcommand{\dirname}[1]{\texttt{\seqsplit{#1}}}
\hypersetup{colorlinks=true,linkcolor=blue,urlcolor=blue}

\title{exp-fips203-mlkem-pke-alg13-16171820-2s1e-k4\\FIPS 203 Alg.13 行 16–20：2s1e MIX + UB 融合\\（FIPS 合规 $\mathbf{s}$/$\mathbf{e}$ 采样版）}
\author{ascendc 工程（预研 customspec）}
\date{2026-06-19}

\begin{document}
\maketitle

\begin{abstract}
本文档为 \dirname{examples/incubating/ml-kem/ml-kem-1024/exp-fips203-mlkem-pke-alg13-16171820-2s1e-k4/} 的\textbf{实现规格说明书}（customspec）。

\textbf{算子语义}：ML-KEM-1024（$k{=}4$，$n{=}256$，$q{=}3329$）K-PKE.KeyGen（FIPS 203 \textbf{Algorithm 13}）行 \textbf{16–20}：
NTT($\mathbf{s}$)、NTT($\mathbf{e}$) $\rightarrow$ $\hat{t}\leftarrow\widehat{A}\circ\hat{\mathbf{s}}+\hat{\mathbf{e}}$ $\rightarrow$ ByteEncode$_{12}$($\hat{t}$/$\hat{\mathbf{s}}$)。

\textbf{工程壳}以已调通探针
\dirname{ascendc-tests/ml-kem/ml-kem-1024/pass-fix-f203-2s1e-alg13-16171820-vec-k4-v2/} 为\textbf{唯一 AscendC 实现参照}（平面 \texttt{mat\_c}、单 TPipe UB 融合、\texttt{stageHatDotOnly} j$\rightarrow$p）。

\textbf{与探针的关键差异（强制）}：本 \texttt{exp} 的 Host \texttt{gen\_data} \textbf{不得}复用探针 shortcuts（\texttt{FIXED\_POLY}$\times 4$、$(e{+}p)\bmod q$ 人工变体）。
$\mathbf{s}{=}(s_0,\ldots,s_{k-1})$ 与 $\mathbf{e}{=}(e_0,\ldots,e_{k-1})$ 须按 FIPS 203 \textbf{Algorithm 13} 行 7–8 + \textbf{Algorithm 8} \texttt{SamplePolyCBD}$_\eta$（ML-KEM-1024：$\eta_1{=}\eta_2{=}2$）\textbf{独立模拟生成}，共 8 条\textbf{互不相同}（高概率）且分布合规的时域多项式。

\textbf{计划 AI Core}：\texttt{KERNEL\_TYPE\_MIX\_AIC\_1\_2}；\texttt{blockDim=1}（\texttt{--aicore 1}）；首版不测多核。

\textbf{阶段}：customspec 闭环；\textbf{禁止}在本目录写可执行实现，直至用户确认规格并明示「可以写代码」。
\end{abstract}

\tableofcontents
\newpage

\section{元数据}

\begin{longtable}{@{}L{3.2cm}L{10.8cm}@{}}
\toprule
项 & 内容 \\
\midrule
用例路径 & \dirname{examples/incubating/ml-kem/ml-kem-1024/exp-fips203-mlkem-pke-alg13-16171820-2s1e-k4/} \\
规格书 & 本文件 \texttt{*-customspec.tex} / \texttt{.pdf} \\
实现参照 & \dirname{ascendc-tests/ml-kem/ml-kem-1024/pass-fix-f203-2s1e-alg13-16171820-vec-k4-v2/}（\texttt{IMPLEMENTATION\_REFERENCE.md}） \\
原理文档 & \dirname{docs/notes/F203-2s1e-NTT内积UB融合技术总结.md} \\
ByteEncode 原理 & \dirname{docs/notes/F203-ByteEncode12-prefetch技术总结.md} \\
探针（encode-only） & \dirname{ascendc-tests/ml-kem/ml-kem-1024/pass-fix-f203-2s1e-byteencode12-vec-k4/} \\
标准 & \dirname{library/documents/NIST.FIPS.203.pdf} Alg.13、Alg.8（CBD） \\
平台 & CANN 9.0.0；Atlas A2（\texttt{Ascend910B4}） \\
核规模 & \texttt{blockDim=1}；1$\times$AIC + 2$\times$AIV / AI Core \\
日期 & 2026-06-19 \\
\bottomrule
\end{longtable}

\section{定位与边界}

\subsection{用例性质}
\begin{itemize}
  \item \textbf{预研 \texttt{exp-*}}：验收 \textbf{仅 I/O 对拍}（\texttt{max\_abs\_diff=0}）。
  \item \textbf{算法片段}：Alg.13 行 16–20（设备侧 MIX 核）；\textbf{不含}行 1–15 的 $\widehat{A}$ 扩展矩阵乘、SHA3 \texttt{G}/\texttt{H}、完整 ML-KEM.KeyGen（Alg.16）。
  \item \textbf{采样在 Host（强制）}：$\mathbf{s}$/$\mathbf{e}$ 由 \texttt{scripts/gen\_data.py}（Python）按 FIPS 203 生成并写入 \texttt{input/src.bin}；设备\textbf{只读} GM \texttt{src}，\textbf{不}在 AscendC 上实时生成 CBD/PRF。
  \item \textbf{常量输入}：$\widehat{A}$ 仍为 Host 提供的 \texttt{a\_hat [16,256]}（与探针相同：固定 \texttt{SEED\_AHAT} 文档化）；后续可换为 Alg.13 行 9–15 全链路。
\end{itemize}

\subsection{与探针 \texttt{vec-k4-v2} 的差异}
\begin{table}[h]
  \centering
  \begin{tabular}{@{}L{3.8cm}L{4.8cm}L{4.8cm}@{}}
    \toprule
    项 & 探针 v2（研究） & 本 \texttt{exp}（交付向） \\
    \midrule
    $\mathbf{s}$ 行 0..3 & 同一 \texttt{FIXED\_POLY} 复制 $\times 4$ & $s_0,\ldots,s_3$ 各自 \texttt{SamplePolyCBD}$_2$ \\
    $\mathbf{e}$ 行 4..7 & 基 $e$ 上 $(e{+}p)\bmod q$ & $e_0,\ldots,e_3$ 各自 \texttt{SamplePolyCBD}$_2$ \\
    分布 & 便于 NTT 手算对拍 & FIPS 203 合规 CBD \\
    行 18 $\hat{e}[p]$ & 每 $p$ 不同（人工） & 每 $p$ 不同（标准采样） \\
    AscendC 壳 & 同左栏 & \textbf{照抄 v2}（平面、UB 融合） \\
    \bottomrule
  \end{tabular}
\end{table}

\subsection{本阶段不做}
\begin{itemize}
  \item \textbf{设备实时生成} $\mathbf{s}$/$\mathbf{e}$（AscendC PRF/CBD、双 launch 预采样）——归属 \dirname{ascendc-tests/} 探针链，\textbf{非}本 \texttt{exp}；Host \texttt{gen\_data} 用 Python（可对照 \texttt{liboqs}）。
  \item Encaps/Decaps、ML-KEM.KeyGen（Alg.16）、FIPS 204。
  \item 多 \texttt{blockDim} 性能调优（须另开 spec 修订）。
  \item 竖堆 \texttt{mat\_c}、\texttt{Matmul<>} 融合、\texttt{SHAT\_PEER}、frozen v1 half-basemul 路线。
\end{itemize}

\section{FIPS 203 输入采样（Host，强制）}
\label{sec:sampling}

\subsection{参数（ML-KEM-1024）}
\begin{itemize}
  \item 模块秩 $k{=}4$；$n{=}256$；$q{=}3329$。
  \item 噪声参数 $\eta_1{=}\eta_2{=}2$（FIPS 203 表 2）。
\end{itemize}

\subsection{语义（Algorithm 13 片段 + Algorithm 8）}
Host 在生成 \texttt{src [8,256]} 时，须等价于下列过程（$d$ 为 32 字节 derand，\texttt{SEED\_D} 文档化，默认建议 \texttt{20260619}）：

\begin{enumerate}
  \item $(\rho,\sigma)\leftarrow \texttt{G}(d\,\|\,k)$（Alg.13 行 5；$k$ 按标准编码为单字节）。
  \item $N\leftarrow 0$。
  \item \textbf{对 $i=0,\ldots,k{-}1$}：$s_i \leftarrow \texttt{SamplePolyCBD}_{\eta_1}(\texttt{PRF}_{\eta_1}(\sigma,N))$；$N\leftarrow N{+}1$。（Alg.13 行 7）
  \item \textbf{对 $i=0,\ldots,k{-}1$}：$e_i \leftarrow \texttt{SamplePolyCBD}_{\eta_2}(\texttt{PRF}_{\eta_2}(\sigma,N))$；$N\leftarrow N{+}1$。（Alg.13 行 8）
\end{enumerate}

\texttt{SamplePolyCBD}$_2$（Alg.8）：由均匀随机字节生成系数，每个系数为两个比特之和减去两个比特之和（中心二项，支撑于 $\{-2,\ldots,2\}$），表示为 $Z_q$ 非负规范整数（与 \texttt{liboqs}/\texttt{mlkem-native} \texttt{mlk\_poly\_cbd2} 一致）。

\subsection{\texttt{src.bin} 布局（与 2s1e 设备契约一致）}
\begin{table}[h]
  \centering
  \begin{tabular}{@{}lll@{}}
    \toprule
    行 & 内容 & 要求 \\
    \midrule
    0..3 & $s_0,\ldots,s_3$ & 四条\textbf{独立} CBD 采样；\textbf{禁止} \texttt{np.tile} 同一行 \\
    4..7 & $e_0,\ldots,e_3$ & 四条\textbf{独立} CBD 采样；\textbf{禁止} $(e{+}p)\bmod q$ 人工变体 \\
    \bottomrule
  \end{tabular}
\end{table}

设备 Stage1 仍将 $\mathbf{s}$ 复制到双 AIV 行块；NTT 后 $\hat{s}_0,\ldots,\hat{s}_3$ \textbf{一般互不相同}（不再要求 $\hat{s}_0{=}\cdots{=}\hat{s}_3$）。

\subsection{参考实现与交叉验证（\texttt{gen\_data} 登记）}
\begin{longtable}{@{}L{4.5cm}L{9.5cm}@{}}
\toprule
来源 & 用途 \\
\midrule
\dirname{thirdparty/liboqs/} ML-KEM-1024 & 黑盒 oracle：相同 $d$ 下 $s_i,e_i$ 系数对照（推荐） \\
\dirname{thirdparty/liboqs/.../sampling.c} & \texttt{mlk\_poly\_cbd2} 语义 \\
\dirname{ascendc-tests/.../vec-k4-v2/scripts/gen\_data.py} & \textbf{仅借鉴} NTT golden 链；\textbf{禁止} \texttt{build\_src\_se()} 采样逻辑 \\
\bottomrule
\end{longtable}

\texttt{gen\_data} 须打印：\texttt{SEED\_D}、8 行唯一性检查（\texttt{assert} 8 行两两不等，或记录 \texttt{max\_abs\_diff} 供审计）。

\section{FIPS 203 Algorithm 13 行 16–20 设备语义}

\subsection{行 16–17（NTT）}
对 $i{=}0,\ldots,k{-}1$：$\hat{s}_i\leftarrow \mathrm{NTT}(s_i)$，$\hat{e}_i\leftarrow \mathrm{NTT}(e_i)$。
本仓以 Tag5T \textbf{三段式}实现（非内核蝶形 \texttt{MlkemNtt}）：
Stage1 limb6 $\rightarrow$ Stage2 int8 MMAD $\times 4$ $\rightarrow$ Stage3 平面 merge + \texttt{stage31\_mod}。

\subsection{行 18（NTT 域内积 + 噪声）}
对每个 $p\in\{0,\ldots,k{-}1\}$：
\[
  \hat{t}[p] = \mathrm{mod}_q\!\left(\sum_{j=0}^{k-1} \widehat{A}[p,j] \circ \hat{s}[j] + \hat{e}[p]\right)
\]
其中 $\circ$ 为 NTT 域多项式乘（Alg.11 \texttt{MultiplyNTTs}）；$j$ 循环内 \textbf{不} final mod；$\hat{e}$ 以向量 \texttt{Add} 累加后 \textbf{一次} \texttt{mod\_q\_final\_vec}（与 v2 \texttt{stageHatDotOnly} 一致）。

\subsection{行 19–20（ByteEncode$_{12}$）}
\label{sec:byteencode}

\textbf{数学（FIPS 203 Alg.5）}：对每个系数取低 12 bit，按标准交织为每 poly \textbf{384 B}；$k{=}4$ 时
$\texttt{ek}\leftarrow \texttt{ByteEncode}_{12}(\hat{t})$、$\texttt{sk}\leftarrow \texttt{ByteEncode}_{12}(\hat{\mathbf{s}})$，
各 \textbf{1536 B}（$4{\times}384$）。

\textbf{设备入口}：\texttt{Aiv2s1eUbPipeline::stageEncodeOut}（与 v2 相同）；对 $p\in\{0,\ldots,k{-}1\}$：
\begin{enumerate}
  \item $\hat{t}[p]$（UB \texttt{ub\_that}）$\rightarrow$ \texttt{poly\_byte\_encode12\_local} $\rightarrow$ GM \texttt{ek} 切片；
  \item $\hat{s}[p]$（UB \texttt{ub\_ntt}，行索引 \texttt{s\_row(p)=p}）$\rightarrow$ 同上 $\rightarrow$ GM \texttt{sk} 切片。
\end{enumerate}

\textbf{双 AIV 分片}（与行 18 一致，\textbf{禁止}单 AIV 写满 polyvec）：
\begin{itemize}
  \item AIV0：$p\in\{0,1\}$，各编 1 poly $\hat{t}$ + 1 poly $\hat{s}$ $\rightarrow$ 本地 768 B + GM 偏移 $p{\cdot}384$；
  \item AIV1：$p\in\{2,3\}$，同上。
\end{itemize}
每 AIV UB 仍握完整 $\hat{s}[0..3]$（Stage1 复制），但 \texttt{sk} 只写本核 $p$ 区间。

\textbf{实现路径（2026-06-19 与 v2 对齐）}：默认 \texttt{BYTE\_ENCODE12\_PREFETCH=1}（整 poly prefetch），
\textbf{非}旧版 tile=32 循环。分发 \texttt{poly\_byte\_encode12\_local}：
\begin{table}[h]
  \centering
  \small
  \begin{tabular}{@{}L{3.6cm}L{8.4cm}@{}}
    \toprule
    宏 & 路径 \\
    \midrule
    \texttt{VEC=1}, \texttt{PREFETCH=1}（\textbf{默认}） & \texttt{poly\_byte\_encode12\_prefetch\_local}：GM ROM Gather 索引 $\rightarrow$ UB；128-wide Gather$\times$2 解交错；\texttt{compute\_b012}；384 B \texttt{DataCopy} \\
    \texttt{VEC=1}, \texttt{PREFETCH=0} & \texttt{poly\_byte\_encode12\_vec\_local}：tile=32，每 tile \texttt{CreateVecIndex}+Gather（legacy，仅对照） \\
    \texttt{VEC=0} & 标量 \texttt{poly\_byte\_encode12\_scalar} \\
    \bottomrule
  \end{tabular}
\end{table}

\textbf{关键文件（fork 自 v2 / \dirname{byteencode12-vec-k4} 探针）}：
\texttt{byte\_encode12\_config.hpp}、\texttt{byte\_encode12\_pair.hpp}、\texttt{byte\_encode12\_vec.hpp}、
\texttt{byte\_encode12\_rom\_tables.*}、\texttt{byte\_encode12\_ub\_load.hpp}；
内核 \texttt{mmad\_custom.cpp} 在 AIV 编译单元 \texttt{\#include byte\_encode12\_rom\_tables.cpp}（与 Alg.11 ROM 相同模式，避免 MIX 双份符号）。

\textbf{UB scratch}：\texttt{HAT\_LINE18\_FULLPOLY=1} 时 encode 工作区在 dot scratch 末尾；
\texttt{PREFETCH=1} 须 \texttt{kPrefetchScratchBytes=5504}（高于 tile 路径 1664 B）。
\texttt{stageEncodeOut} 循环体与 v2 一致，仅 scratch 预算与 \texttt{encode\_ws} 槽位绑定变更。

\textbf{Golden}：Host \texttt{byte\_encode12\_ref}；I/O 与 v2 \texttt{mixPass=7} preset 字节级一致（见 encode-only 探针 \dirname{byteencode12-vec-k4}）。

\textbf{性能参考（SIM，910B4，非门禁）}：v2 全链路 prefetch \textbf{$\sim$77958 tick}（encode 边际 $\sim$+12421 vs 无 encode）；
encode-only 探针 prefetch \textbf{$\sim$17429 tick}。本 \texttt{exp} fork 同步 v2 后 tick 应同量级（FIPS CBD 输入不改变设备 encode 路径）。

\section{计算语义：三段式 NTT + 平面 \texttt{mat\_c}}

\subsection{流水线}
\begin{enumerate}
  \item \textbf{Stage1}（AIV \texttt{Aiv2s1eSplit}）：\texttt{src [8,256]} $\rightarrow$ \texttt{S0} limb6 int8；双 $\hat{\mathbf{s}}$ 行块 + $\hat{\mathbf{e}}$ 对半。
  \item \textbf{Stage2}（AIC \texttt{AicMmad}$\times 4$）：even/odd $\times$ lo/hi $\rightarrow$ 四路临时 \texttt{int32}。
  \item \textbf{Stage2 后 pack}（AIV \texttt{Aiv2s1ePackMatCPlanar}）：$\rightarrow$ \texttt{mat\_c\_planar [96,128]}。
  \item \textbf{Stage3 + 行18 + 行19–20}（AIV \texttt{Aiv2s1eUbPipeline}）：平面 merge $\rightarrow$ $\hat{\mathbf{s}}$/$\hat{\mathbf{e}}$ UB；j$\rightarrow$p 内积 $\rightarrow$ $\hat{t}$ UB；ByteEncode $\rightarrow$ GM \texttt{ek}/\texttt{sk}。
\end{enumerate}

\subsection{平面行号（禁止混用）}
\[
  \texttt{planar\_row}(\mathit{slot},\mathit{limb},\mathit{half})
  = \mathit{half}\cdot 48 + \mathit{slot}\cdot 4 + \mathit{limb}
\]
slot：0..3 $\hat{\mathbf{s}}$ AIV0；4..7 $\hat{\mathbf{s}}$ AIV1（同值复制）；8..9 $\hat{\mathbf{e}}$ AIV0；10..11 $\hat{\mathbf{e}}$ AIV1。

\textbf{NTT S1–S3 禁止 Gather}；仅 bulk \texttt{DataCopy} 四 limb 行。

\section{数据搬运与 UB 不变量}

\subsection{驻留 vs 落盘}
\begin{itemize}
  \item \textbf{驻留}：S3 后 $\hat{\mathbf{s}}$、$\hat{\mathbf{e}}$、$\hat{t}$ 在 UB（\texttt{ubNttBuf\_}/\texttt{ubThatBuf\_}）直至 ByteEncode 结束（§\ref{sec:byteencode}）。
  \item \textbf{ByteEncode 输入}：$\hat{t}[p]$、$\hat{s}[p]$ 已在 UB；\textbf{禁止} encode 阶段再从 GM 读 \texttt{dst}/\texttt{t\_hat} preset（全链路 \texttt{mixPass=0}）。
  \item \textbf{落盘}：对拍 \texttt{dst}/\texttt{t\_hat} dump 仅在阶段结束；\texttt{ek}/\texttt{sk} 为算法 GM 输出。
  \item \textbf{禁止}：S3 后 $\hat{\mathbf{s}}$ GM 绕路；\textbf{禁止} \texttt{SHAT\_PEER}；\textbf{禁止}嵌套第二 \texttt{TPipe}。
\end{itemize}

\subsection{行 18 数据面}
\begin{itemize}
  \item $\widehat{A}[p,j]$：只读 GM \texttt{a\_hat}（行主序 $(p{\cdot}K{+}j){\cdot}N{+}c$）。
  \item $\hat{s}[j]$：只读 UB \texttt{ub\_ntt}（每 AIV 握完整 $j\in[0,3]$）。
  \item Alg.11 ROM（$\gamma$/Gather/interleave）：Init \texttt{DataCopy} 进 UB，Compute 禁止 \texttt{SetValue} 重填。
\end{itemize}

\section{MIX 调度与 \texttt{mixPass}}

\subsection{FSM（\texttt{blockDim=1}）}
\begin{enumerate}
  \item \texttt{AIV\_SPLIT}：Stage1；末 \texttt{SET}。
  \item \texttt{AIC\_MMAD}：\texttt{WAIT}；MMAD$\times 4$；末 \texttt{SET}。
  \item \texttt{AIV\_PACK}：\texttt{WAIT}；平面 pack。
  \item \texttt{Aiv2s1eUbPipeline}：S3 / 行18 / Encode（无额外核间同步）。
\end{enumerate}

\subsection{\texttt{mixPass} 调试}
与 v2 \texttt{IMPLEMENTATION\_REFERENCE.md} §4 相同：0=全链路；1..7 分阶段；CPU 可 5$\rightarrow$4 checkpoint。

\section{AI Core 规模（强制专节）}
\label{sec:aicore}

\begin{longtable}{@{}L{2.8cm}L{10.5cm}@{}}
\toprule
项 & 本用例约定 \\
\midrule
核类型 & \texttt{KERNEL\_TYPE\_MIX\_AIC\_1\_2}（1 Cube + 2 Vector / 1 AI Core） \\
计划 AI Core 数 & \textbf{仅 1}：\texttt{blockDim=1}；\texttt{run.sh --aicore 1} \\
数据划分 & AIV0：$p\in\{0,1\}$，ê 行 4..5；AIV1：$p\in\{2,3\}$，ê 行 6..7；双核各握完整 $\hat{s}[0..3]$ \\
选型理由 & 与 v2 探针已验证路径一致；单趟 launch 完成行 16–20；多核须另开 spec \\
多档验证 & 首版\textbf{不}测；晋级 stable 前可增 \texttt{aicore} 档位 \\
\bottomrule
\end{longtable}

\section{I/O 与 Golden 契约}

\subsection{文件（\texttt{gen\_data} 与 \texttt{main} 逐字一致）}
\begin{longtable}{@{}L{3.8cm}lll@{}}
\toprule
路径 & 形状 & dtype & 语义 \\
\midrule
\texttt{input/src.bin} & $[8,256]$ & int32 & §\ref{sec:sampling} FIPS CBD $\mathbf{s}\|\mathbf{e}$ \\
\texttt{input/a\_hat.bin} & $[16,256]$ & int32 & $\widehat{A}$；\texttt{SEED\_AHAT=20260615} \\
\texttt{input/lut\_even/odd\_stacked.bin} & $[512,128]$ & int8 & Tag5T LUT（\texttt{transpose\_mlkem\_luts\_i8.h}） \\
\texttt{input/tiling.bin} & $(256,4,\mathit{mixPass})$ & int32 & 与 v2 \texttt{tiling.h} \\
\texttt{output/dst.bin} & $[12,256]$ & int32 & NTT 后 $\hat{\mathbf{s}}$/$\hat{\mathbf{e}}$ \\
\texttt{output/t\_hat.bin} & $[4,256]$ & int32 & 行 18 \\
\texttt{output/ek\_polyvec.bin} & $4{\times}384$ & uint8 & 行 19 \\
\texttt{output/sk\_polyvec.bin} & $4{\times}384$ & uint8 & 行 20 \\
\texttt{output/golden\_*.bin} & 同上 & 同上 & Host 参考 \\
\bottomrule
\end{longtable}

\subsection{Golden 生成链（NTT 与 v2 同构）}
\begin{verbatim}
src (FIPS CBD s||e)
  -> encode_2s1e_s0 -> mat_c_tmp_golden -> pack_mat_c_planar
  -> golden_dst_from_planar (merge + stage31_mod)
  -> hat_inner_product_ref (dot / add)
  -> byte_encode12_ref (ek/sk)
\end{verbatim}
\textbf{禁止}在 \texttt{gen\_data} 内对 NTT 逐行调用 C \texttt{MlkemNtt()}；须与设备三段式同构（与 v2 §2.1 相同）。

\subsection{Golden 与设备 mod 解耦}
行 18 golden：\texttt{HAT\_GOLDEN\_MOD\_VARIANT=0}（int64 floor）；设备：\texttt{F203\_MOD\_VARIANT=1} Barrett 向量。对拍比最终 \texttt{int32} 多项式。

\subsection{验收}
\begin{verbatim}
bash run.sh -r cpu -v Ascend910B4
bash run.sh -r sim -v Ascend910B4
python3 scripts/verify_result.py
\end{verbatim}
\textbf{默认}即全链路（\texttt{HAT\_LINE18\_DOT\_ONLY=0}、\texttt{HAT\_BYTE\_ENCODE=1}、\texttt{BYTE\_ENCODE12\_PREFETCH=1}，见 \texttt{run.sh}）；调试分段须显式覆盖。
要求 \texttt{max\_abs\_diff=0}。SIM 须 \texttt{source scripts/camodel\_sim\_log.sh}。

\section{AscendC 实现规划}

\subsection{参照 fork 顺序}
\begin{enumerate}
  \item \textbf{壳}：自 \dirname{ascendc-tests/ml-kem/ml-kem-1024/pass-fix-f203-2s1e-alg13-16171820-vec-k4-v2/} 复制
    \texttt{mmad\_custom.cpp}、\texttt{2s1e\_post\_ntt\_ub.hpp}、\texttt{aiv\_func.hpp}、\texttt{tiling.h}、
    \texttt{alg11\_*}、\texttt{hat\_*}、\texttt{byte\_encode12\_*}（含 \texttt{rom\_tables.*}、\texttt{ub\_load.hpp}）、\texttt{cmake/}、\texttt{run.sh}。
  \item \textbf{ByteEncode 宏}：默认 \texttt{BYTE\_ENCODE12\_VEC=1}、\texttt{BYTE\_ENCODE12\_SCATTER\_VEC=1}、\texttt{BYTE\_ENCODE12\_PREFETCH=1}（与 v2 \texttt{run.sh} 一致）。
  \item \textbf{改} \texttt{scripts/gen\_data.py}：实现 §\ref{sec:sampling}；删除 \texttt{FIXED\_POLY}/\texttt{NTTS2S1E\_E\_POLY\_SEED} 探针逻辑。
  \item \textbf{改} \texttt{verify\_result.py}：保留全链路对拍；\textbf{移除}「$\hat{s}$ 四行相同」探针专用断言（或改为可选 \texttt{--probe-layout}）。
  \item 内核文件名建议：\texttt{f203\_alg13\_16171820\_2s1e\_custom.cpp}（与探针 \texttt{mmad\_custom} 区分）。
\end{enumerate}

\subsection{计划目录（实现阶段）}
\begin{verbatim}
exp-fips203-mlkem-pke-alg13-16171820-2s1e-k4/
  *-实现方案-customspec.tex / .pdf
  STATUS.md
  f203_alg13_16171820_2s1e_custom.cpp
  main.cpp, CMakeLists.txt, run.sh
  scripts/gen_data.py, verify_result.py
  (headers forked from vec-k4-v2)
\end{verbatim}

\subsection{分步验收}
\begin{table}[h]
  \centering
  \begin{tabular}{@{}cll@{}}
    \toprule
    步 & 内容 & 验收 \\
    \midrule
    S0 & 本 customspec PDF & 文档就绪 \\
    S1 & \texttt{gen\_data}：FIPS CBD $\mathbf{s}$/$\mathbf{e}$ + golden 链 & liboqs 系数对照；8 行互异 \\
    S2 & fork v2 内核 + 新 \texttt{src} & CPU \texttt{max\_diff=0} \\
    S3 & SIM 全链路 & tick 参考 $<10^5$；v2 prefetch 基准 $\sim$77958（附录，非门禁） \\
    S4 & \texttt{npu} & 有实机时 \\
    \bottomrule
  \end{tabular}
\end{table}

\section{全量 API 清单（主路径，CANN 9.0.0）}
\label{sec:api}

说明：分类按 PDF 第 2 章；在线页为 \texttt{atlasascendc\_api\_07\_*}。

\begin{longtable}{@{}L{2.4cm}L{1.4cm}L{2.2cm}L{0.7cm}L{3.2cm}L{3.5cm}@{}}
\toprule
API & 分类 & 在线 ID & A2 & 本用例 dtype & 备注 \\
\midrule
\endhead
\apiname{KERNEL\_TASK\_TYPE\_DEFAULT} & Utils & 编程指南 & $\checkmark$ & \texttt{MIX\_AIC\_1\_2} & 融合核 \\
\apiname{CrossCoreSetFlag} & 同步 & 07 列表 & $\checkmark$ & event id & AIV$\rightarrow$AIC \\
\apiname{CrossCoreWaitFlag} & 同步 & 07 列表 & $\checkmark$ & event id & AIC$\leftarrow$AIV \\
\apiname{PipeBarrier} & 同步 & 2.3 & $\checkmark$ & \texttt{PIPE\_ALL}/\texttt{MTE2} & \texttt{KYBER\_PIPE\_ALL} \\
\apiname{DataCopy} & 2.3.1 & 07\_0101 & $\checkmark$ & int32/int8 GM$\leftrightarrow$UB & 平面四 limb；\textbf{无 Gather} \\
\apiname{ShiftRight} & 2.3.3 & 07\_0059 & $\checkmark$ & int32 & Stage1 $hi$ \\
\apiname{Muls} & 2.3.3 & 07\_0055 & $\checkmark$ & int32 & Stage1 $hi{\cdot}64$ \\
\apiname{Sub} & 2.3.3 & 07\_0036 & $\checkmark$ & int32 & Stage1 $lo$ \\
\apiname{Add} & 2.3.3 & 07\_0035 & $\checkmark$ & int32 & 行18 lazy $\Sigma$；$+\hat{e}$ \\
\apiname{Cast} & 2.3.3 & 07\_0073 & $\checkmark$ & i32$\rightarrow$i16$\rightarrow$half$\rightarrow$i8 & Stage1 窄化链 \\
\apiname{Gather} & 2.3.3 & 07 列表 & $\checkmark$ & int32 索引 & Alg.11 basemul；\textbf{ByteEncode} 解交错（§\ref{sec:byteencode}）；NTT S1–S3 禁用 \\
\apiname{TPipe}/\apiname{TQue}/\apiname{TBuf} & 2.3 & 07\_0108 & $\checkmark$ & UB 资源 & 单 slim TPipe \\
\apiname{AicMmad} & Cube & merged\_kyber & $\checkmark$ & int8$\times$int8$\rightarrow$int32 & Stage2 $\times 4$ \\
\bottomrule
\end{longtable}

\subsection{评估过但未采用}
\begin{longtable}{@{}L{2.6cm}L{2.2cm}L{4.8cm}@{}}
\toprule
API/路线 & 结论 \\
\midrule
\apiname{Matmul<>} / \texttt{REGIST\_MATMUL\_OBJ} & NTT 融合已冻结 \\
\apiname{SHAT\_PEER} GM 交换 & 2s1e 每核握完整 $\hat{\mathbf{s}}$ \\
竖堆 \texttt{mat\_c} + Gather S3 & 平面 bulk 取代 \\
探针 \texttt{FIXED\_POLY} 输入 & 本 \texttt{exp} 禁止 \\
\bottomrule
\end{longtable}

\section{逐步覆盖率评估}

\begin{longtable}{@{}L{1.2cm}L{4.0cm}L{3.8cm}L{2.6cm}@{}}
\toprule
段 & 数学/操作 & API & 覆盖 \\
\midrule
S1 & limb6 拆分 & \texttt{ShiftRight}/\texttt{Muls}/\texttt{Sub}/\texttt{Cast} & 100\% 向量 \\
S2 & int8 MMAD & \texttt{AicMmad} & Cube \\
S3 & merge+mod & \texttt{DataCopy}+RouteA 向量 & 100\%（无 Gather） \\
18 & Alg.11$\times k$ + $\Sigma$ + $\hat{e}$ + mod & \texttt{compute\_on\_ub}/\texttt{Add}/\texttt{mod\_q\_final\_vec} & 100\% 向量 \\
19–20 & ByteEncode$_{12}$ & \texttt{poly\_byte\_encode12\_prefetch\_local}（默认） & 100\% 向量；ROM+128 Gather；双 AIV 各 2 poly \\
\midrule
\multicolumn{3}{l}{Host CBD 采样} & \textbf{Host Python}（\textbf{非} AI Core） \\
\bottomrule
\end{longtable}

\textbf{标量缺口}：\texttt{gen\_data} 中 FIPS \texttt{G}/\texttt{PRF}/CBD；CPU\_DEBUG 下行 18 标量回退（与 v2 相同）。

\section{附录：SIM tick 参考（910B4，非门禁）}

\begin{table}[h]
  \centering
  \small
  \begin{tabular}{@{}L{5.5cm}r@{}}
    \toprule
    配置（v2 / 同步后 exp） & Total tick \\
    \midrule
    全链路 + \texttt{BYTE\_ENCODE12\_PREFETCH=1}（默认） & $\sim$77958 \\
    全链路 + \texttt{PREFETCH=0}（tile32 legacy） & $\sim$85991 \\
    encode-only 探针 prefetch & $\sim$17429 \\
    \bottomrule
  \end{tabular}
\end{table}
FIPS CBD 输入仅改变 Host \texttt{src}；设备 encode 路径与探针 v2 相同。本 \texttt{exp} 在 2026-06-19 前 fork 的 tick $\sim$86111 为 \textbf{tile32 时代}测量；同步 prefetch 后应降至 $\sim$77958 量级。



\section{已废弃路径}

\begin{itemize}
  \item \dirname{ascendc-tests/frozen/frozen-fix-f203-2s1e-alg13-16171820-vec-k4/} half-basemul 路线。
  \item v2 \texttt{build\_src\_se()} 的 \texttt{FIXED\_POLY} 与 $(e{+}p)$ 变体——\textbf{本 exp 明确禁止}。
\end{itemize}

\section{关联文档}

\begin{itemize}
  \item 探针：\dirname{ascendc-tests/ml-kem/ml-kem-1024/pass-fix-f203-2s1e-alg13-16171820-vec-k4-v2/IMPLEMENTATION\_REFERENCE.md}
  \item 原理：\dirname{docs/notes/F203-2s1e-NTT内积UB融合技术总结.md}
  \item ByteEncode：\dirname{docs/notes/F203-ByteEncode12-prefetch技术总结.md}；\dirname{qa/2026-06/2026-06-19-ByteEncode12-only探针与prefetch实验.md}
  \item 讨论：\dirname{qa/2026-06/2026-06-18-内积布局与NTT内积UB融合讨论.md}
  \item 标准：\dirname{library/documents/NIST.FIPS.203.pdf}
  \item liboqs 对照：\dirname{thirdparty/liboqs/}
\end{itemize}

\end{document}
